\chapter{Complete Ricci flows of bounded curvature}
\label{chap:complete-ricci-flows}
\label{newcomp2}

We work with complete Ricci flows whose curvature is bounded on compact
time intervals. The backward time variable is $\tau=T-t$.

\begin{definition}[Complete Ricci flow of bounded curvature]
\label{def:complete-bounded-curvature-flow}
\uses{def:ricci-flow}

Let $(M,g(t)),\ a\le t\le b$, be a  Ricci flow. We say that the flow
is {\em complete of bounded curvature} if for each $t\in [a,b]$ the
Riemannian manifold $(M,g(t))$ is complete and if there is
$C<\infty$ such that $|\Rm|(p,t)\le C$ for all $p\in M$ and all
$t\in [a,b]$. Let $I$ be an interval and let $(M,g(t)),\ t\in I$, be
a Ricci flow. Then we say that the flow is {\em complete with
curvature locally bounded in time} if for each compact subinterval
$J\subset I$ the restriction of the flow to $(M,g(t)),\ t\in J$, is
complete of bounded curvature.
\end{definition}




\section{The functions $L_{x}$ and $l_{x}$}

Throughout this chapter $(M,g(t))$, $0\le t\le T<\infty$, is a Ricci
flow and $\tau=T-t$.

\subsection{Existence of ${\mathcal L}$-geodesics}


We assume here that $(M,g(t)),\ 0\le t\le T<\infty$, is a Ricci flow
which is complete of bounded curvature. In Shi's Theorem
(Theorem~\ref{shi}) we take $K$ equal to the bound of the norm of
the Riemannian curvature on $M\times [0,T]$, we take $\alpha=1$, and
we take $t_0=T$. It follows from Theorem~\ref{shi} that there is a
constant $C(K,T)$ such that $|\nabla R(x,t)|\le C/t^{1/2}$. Thus,
for any $\epsilon>0$ we have a uniform bound for $|\nabla R|$ on
$M\times [\epsilon,T]$. Also, because of the uniform bound for the
Riemann curvature and the fact that $T<\infty$, there is a constant
$C$, depending on the curvature bound and $T$ such that
\begin{equation}\label{compmetric}
C^{-1}g(x,t)\le g(x,0)\le Cg(x,t)
\end{equation}
for all $(x,t)\in M\times [0,T]$.





\begin{lemma}[Existence of minimizing L-geodesics]
\label{lem:l-geodesic-existence}
\uses{def:complete-bounded-curvature-flow}
\label{geoexist} Assume that $M$ is connected.
Given $p_1,p_2\in M$ and $0\le\tau_1<\tau_2\le T$, there is
 a minimizing ${\mathcal L}$-geodesic:
$\gamma\colon[\tau_{1},\tau_{2}]\to M\times [0,T]$ connecting
$(p_1,\tau_1)$ to $(p_2,\tau_2)$.
\end{lemma}

\begin{proof}\sketch For any curve $\gamma$ parameterized by backward time, we set
$\bar \gamma$ equal to the path in $M$ that is the image under
projection of $\gamma$. We set $A(s)=\bar\gamma'(s)$. Define
$$c((p_1,\tau_1),(p_2,\tau_{2})) = \mathrm{inf}
\{\mathcal{L}(\gamma)|\gamma\colon [\tau_{1},\tau_{2}] \rightarrow
M\times [0,T], \bar\gamma(\tau_{1}) = p_1, \bar\gamma(\tau_{2})=
p_2\}.$$ From Equation~(\ref{seqn}) we see that the infimum exists
since, by assumption, the curvature is uniformly bounded (below).
Furthermore, for a minimizing sequence $\gamma_{i}$, we have
$\int_{s_1}^{s_{2}}\abs{A_i(s)}^{2}ds \leq C_0$, for some constant
$C_0$, where $s_i=\sqrt{\tau_i}$ for $i=1,2$. It follows from this
and the inequality in Equation~(\ref{compmetric}) that there is a
constant $C_1<\infty$ such that for all $i$ we have
$$\int_{s_{1}}^{s_{2}} \abs{
A_{i}}^{2}_{g(0)}d\tau \leq C_1.$$ Therefore the sequence $\{\gamma_{i}\}$ is
uniformly continuous with respect to the metric $g(0)$; by Cauchy-Schwarz we
have
$$\abs{\bar\gamma_{i}(s) - \bar\gamma_{i}(s')}_{g(0)} \leq
\int_{s'}^{s} \abs{A_{i}}_{g(0)}ds \leq \sqrt{C_1}\sqrt{s-s'}.$$ By
the uniform continuity, we see that a subsequence of the $\gamma_i$
converges uniformly pointwise to a continuous curve $\gamma$
parameterized by $s$, the square root backward time. By passing to a
subsequence we can arrange that the $\gamma_i$ converge weakly in
$H^{2,1}$. Of course, the limit in $H^{2,1}$ is represented by the
continuous limit $\gamma$. That is to say, after passing to a
subsequence, the $\gamma_i$ converge uniformly and weakly in
$H^{2,1}$ to a continuous curve $\gamma$. Let $A(s)$ be the
$L^2$-derivative of $\gamma$. Weak convergence in $H^{2,1}$ implies
that $\int_{s'}^s|A(s)|^2ds\le \mathrm{lim}_{i\rightarrow\infty}\int_{s'}^s|A_i(s)|^2ds$, so that
${\mathcal L}(\gamma)\le \mathrm{lim}_{i\rightarrow\infty}{\mathcal
L}(\gamma_i)$. This means that $\gamma$ minimizes the ${\mathcal
L}$-length. Being a minimizer of ${\mathcal L}$-length, $\gamma$
satisfies the Euler-Lagrange equation and is smooth by the
regularity theorem of differential equations. This then is the
required minimizing ${\mathcal L}$-geodesic from $(p_1,\tau_1)$ to
$(p_2,\tau_2)$.
\end{proof}


Let us now show that it is always possible to uniquely extend
${\mathcal L}$-geodesics up to time $T$.

\begin{lemma}[Unique extension of L-geodesics to time T]
\label{lem:l-geodesic-extension}
\uses{def:complete-bounded-curvature-flow, lem:l-geodesic-existence, lem:l-geodesic-min-max-estimate}
 For any $0\le \tau_1<\tau_2<T$ suppose that
 $\gamma\colon [\tau_1,\tau_2]\to M\times [0,T]$ is an ${\mathcal L}$-geodesic.
 Then $\gamma$ extends uniquely to an ${\mathcal L}$-geodesic
$\gamma\colon [0,T)\to M\times (0,T]$.
\end{lemma}


\begin{proof}\sketch
We work with the parameter $s=\sqrt{\tau}$. According to
Equation~(\ref{reparam}), we have
$$\nabla_{\gamma'(s)}\gamma'(s)=2s^2\nabla
R-4s\Ric(\gamma'(s),\cdot).$$ This is an everywhere
non-singular ODE. Since the manifolds $(M,g(t))$ are complete and
their metrics are uniformly related as in
Inequality~(\ref{compmetric}), to show that the solution is defined
on the entire interval $s\in [0,\sqrt{T})$ we need only show that
there is a uniform bound to the length, or equivalently the energy
of $\gamma$ of any compact subinterval of $[0,T)$ on which it is
defined. Fix $\epsilon>0$. It follows immediately from
Lemma~\ref{min-max}, and the fact that the quantities $R$, $|\nabla
R|$ and $|\Rm|$ are bounded on $M\times [\epsilon,T]$, that
there is a bound on $\mathrm{max}|\gamma'(s)|$ in terms of
$|\gamma'(\tau_1)|$, for all $s\in [0,\sqrt{T-\epsilon}]$ for which
$\gamma$ is defined. Since $(M,g(0))$ is complete, this, together
with a standard extension result for second-order ODE's, implies
that $\gamma$ extends uniquely to the entire interval
$[0,\sqrt{T-\epsilon}]$. Changing the variable from $s$ to
$\tau=s^2$ shows that the ${\mathcal L}$-geodesic extends uniquely
to the entire interval $[0,T-\epsilon]$. Since this is true for
every $\epsilon>0$, this completes the proof.
\end{proof}

Let $p\in M$ and set $x=(p,T)\in M\times [0,T]$. Recall that from
Definition~\ref{injdefn} for every $\tau>0$, the injectivity set
$\widetilde{\mathcal U}_{x}(\tau)\subset T_pM$ consists of all $Z\in
T_pM$ for which (i) the ${\mathcal L}$-geodesic
$\gamma_Z|_{[0,\tau]}$ is the unique minimizing ${\mathcal
L}$-geodesic from $x$ to its endpoint, (ii) the differential of
${\mathcal L}\mathrm{exp}_{x}^\tau$ is an isomorphism at $Z$, and (iii)
for all  $Z'$ sufficiently close to
 $Z$ the ${\mathcal L}$-geodesic $\gamma_{Z'}|_{[0,\tau]}$ is the unique minimizing
 ${\mathcal L}$-geodesic to its endpoint. The image
of $\widetilde{\mathcal U}_{x}(\tau)$ is denoted ${\mathcal U}_{x}(\tau)\subset
M$




The existence of minimizing ${\mathcal L}$-geodesics from $x$ to
every point of $M\times (0,T)$ means that the functions $L_{x}$ and
$l_{x}$ are defined on all of $M\times (0,T)$. This leads to:

\begin{definition}[The functions L_x and l_x on complete space-time]
\label{def:reduced-length-global}
\uses{lem:l-geodesic-existence, def:l-injectivity-domain}
 Suppose that $(M,g(t)),\ 0\le t\le T<\infty$, is a Ricci
flow, complete of bounded curvature.  We define the function
$L_{x}\colon M\times [0,T)\to \Ar$ by assigning to each $(q,t)$ the
length of any ${\mathcal L}$-minimizing ${\mathcal L}$-geodesic from
$x$ to $y=(q,t)\in M\times [0,T)$. Clearly, the restriction of this
function to ${\mathcal U}_{x}$ agrees with the smooth function
$L_{x}$ given in Definition~\ref{L_x}. We define
 $L_{x}^\tau\colon M\to \Ar$ to be the restriction of $L_x$ to
 $M\times \{T-\tau\}$.
Of course, the restriction of $L_{x}^\tau$ to ${\mathcal
U}_{x}(\tau)$ agrees with the smooth function $L_{x}^\tau$ defined
in the last chapter. We define $l_{x}\colon M\times [0,T)\to \Ar$ by
$l_{x}(y)=L_{x}(y)/2\sqrt{\tau}$, where, as always $\tau=T-t$, and
we define $l_{x}^\tau(q)=l_{x}(q,T-\tau)$.
\end{definition}





\subsection{Results about $l_{x}$ and ${\mathcal
U}_{x}(\tau)$}

Now we come to our main result about the nature of ${\mathcal
U}_{x}(\tau)$ and the function $l_{x}$
in the context of Ricci flows which are complete and of bounded
curvature.



\begin{proposition}[Lipschitz regularity and diffeomorphism properties of L_x and l_x]
\label{prop:l-exp-lipschitz-complete}
\uses{def:reduced-length-global}
\label{lipcomplete}
Let $(M,g(t)),\ 0\le t\le T<\infty$, be a Ricci flow that is
complete and of bounded curvature. Let $p\in M$, let $x=(p,T)\in
M\times [0,T]$, and let $\tau\in (0,T)$.
\begin{enumerate}
\item[(1)] The functions $L_{x}$ and $l_{x}$ are locally Lipschitz functions
on $M\times (0,T)$.
\item[(2)] ${\mathcal L}\mathrm{exp}_{x}^\tau$
 is a diffeomorphism from $\widetilde{\mathcal U}_{x}(\tau)$ onto
an open subset ${\mathcal U}_{x}(\tau)$ of $M$.
\item[(3)] The complement of ${\mathcal U}_{x}(\tau)$ in $M$ is a closed
subset of $M$ of zero Lebesgue measure. \item[(4)]
 For every
$\tau<\tau'<T$ we have
$${\mathcal U}_{x}(\tau')\subset {\mathcal U}_{x}(\tau).$$
\end{enumerate}
\end{proposition}

\begin{proof}\sketch
\uses{thm:shi-derivative-estimates,prop:reduced-length-lipschitz-precise,prop:l-exponential-diffeomorphism,prop:injectivity-set-monotone,cor:reduced-length-full-measure}
By Shi's Theorem (Theorem~\ref{shi}) the curvature bound on $M\times
[0,T]$ implies that for each $\epsilon>0$ there is a bound for
$|\nabla R|$ on $M\times (\epsilon,T]$. Thus,
Proposition~\ref{lipaty} shows that $L_{x}$ is a
 locally Lipschitz function on $M\times (\epsilon,T)$. Since this is
true for every $\epsilon>0$, $L_x$ is a locally Lipschitz function
on $M\times (0,T)$. Of course, the same is true for $l_{x}$. The
second statement is contained in Proposition~\ref{diffeo}, and the
last one is contained in Proposition~\ref{inclusion}. It remains to
prove the third statement, namely that the complement of ${\mathcal
U}_{x}(\tau)$ is closed nowhere dense. This follows immediately from
Corollary~\ref{fullmeasure} since $|\Ric|$ and $|\nabla R|$ are
bounded on $F=M\times [T-\tau,T]$.
\end{proof}

\begin{corollary}[Continuity and almost-everywhere smoothness of l_x]
\label{cor:l-continuous-ae-smooth}
\uses{prop:l-exp-lipschitz-complete, def:reduced-length-global}

The function $l_x$ is a continuous function on $M\times (0,T)$ and
is smooth on the complement of a closed subset ${\mathcal C}$ that
has the property that its intersection with each $M\times \{t\}$ is
of zero Lebesgue measure in $M\times \{t\}$. For each $\tau\in
(0,T)$ the gradient $\nabla l^\tau_x$ is then a smooth vector field
on the complement of ${\mathcal C}\cap M_{T_\tau}$, and it is a
locally essentially bounded vector field in the following sense. For
each $q\in M$ there is a neighborhood $V\subset M$ of $q$ such that
the restriction of $|\nabla l_x^\tau|$ to $V\setminus\left(V\cap
{\mathcal C}\right)$ is a bounded smooth function. Similarly,
$\partial l_x/\partial t$ is an essentially bounded smooth vector
field on $M\times (0,T)$.
\end{corollary}

\section{A bound for $\mathrm{min}\, l^\tau_x$}

We continue to assume that we have a Ricci flow $(M,g(t)),\ 0\le
t\le T<\infty$, complete and of bounded curvature and a point
$x=(p,T)\in M\times [0,T]$. Our purpose here is to extend the first
differential inequality given in
 Corollary~\ref{lformula} to a differential
inequality in the weak or distributional sense for $l_{x}$ valid on
all of $M\times (0,T)$. We then use this to establish that $\mathrm{min}_{q\in M} l^\tau_{x}(q)\le n/2$ for all $0<\tau<T$.

In establishing inequalities in the non-smooth case the notion of a
support function or a barrier function is often convenient.

\begin{definition}[Upper barrier function]
\label{def:upper-barrier}
Let $P$ be a smooth manifold and let $f\colon P\to \Ar$ be a continuous
function. An {\em upper barrier} for $f$ at $p\in P$ is a
smooth function $\varphi$ defined on a neighborhood of $p$ in $P$, say $U$,
satisfying $\varphi(p)=f(p)$ and $\varphi(u)\ge f(u)$ for all $u\in U$, 
\end{definition}





\begin{proposition}[Differential inequality for l_x in the barrier sense]
\label{prop:l-barrier-inequality}
\uses{def:upper-barrier, def:reduced-length-global}
\label{weaksense}
Let $(M,g(t)),\ 0\le t\le T<\infty$, be an $n$-dimensional Ricci
flow,  complete  of bounded curvature. Fix a point $x=(p,T)\in
M\times [0,T]$, and for any $(q,t)\in M\times [0,T]$, set
$\tau=T-t$. Then for any $(q,t)$, with $0<t<T$, we have
$$\frac{\partial l_x}{\partial \tau}(q,\tau) +\triangle l_x(q,\tau)\le
\frac{(n/2)-l_x(q,\tau)}{\tau}$$ in the barrier sense. This means
that
 for each $\epsilon>0$ there is a neighborhood $U$ of $(q,t)$ in
 $M\times [0,T]$
 and an upper barrier $\varphi$ for $l_x$ at this point defined on $U$
satisfying
$$\frac{\partial \varphi}{\partial \tau}(q,\tau)+\triangle \varphi(q,\tau)
\le \frac{(n/2)-l_x(q,\tau)}{\tau}+\epsilon.$$
\end{proposition}

\begin{proof}\sketch
\uses{lem:l-geodesic-existence,cor:reduced-length-pde-inequalities,prop:l-exponential-diffeomorphism-along-geodesic,lem:reduced-length-differential-inequalities,lem:k-quantity-continuity}
If $(q,T-\tau)\in {\mathcal U}_{x}$, then $l_x$ is smooth near $(q,T-\tau)$,
and the result is immediate from the first inequality in
Corollary~\ref{lformula}.


Now consider a general point $(q,t=T-\tau)$ with $0<t<T$. According to
Lemma~\ref{geoexist} there is  a minimizing ${\mathcal L}$-geodesic $\gamma$
from $x=(p,T)$ to $(q,t=T-\tau)$. Let $\gamma$ be any minimizing ${\mathcal
L}$-geodesic between these points. Fix $0<\tau_1<\tau$ let $q_1=\gamma(\tau_1)$
and set $t_1=T-\tau_1$. Even though $q_1$ is contained in the $t_1$ time-slice,
we keep $\tau=T-t$ so that paths beginning at $q_1$ are parameterized by
intervals in the $\tau$-line of the form $[\tau_1,\tau']$ for some $\tau'<T$.
Consider ${\mathcal L}\mathrm{exp}_{(q_1,t_1)}\colon T_{q_1}M\times (\tau_1,T)\to
M\times (0,t_1)$. According to Proposition~\ref{tausmooth} there is a
neighborhood $\widetilde V$ of $\{\sqrt{\tau_1}X_\gamma(\tau_1)\}\times
(\tau_1,\tau]$ which is mapped diffeomorphically  by ${\mathcal L}\mathrm{exp}_{(q_1,t_1)}\colon T_{q_1}M\times (\tau_1,\tau)\to M\times (0,t_1)$ onto a
neighborhood $V$ of $\gamma((\tau_1,\tau))$. (Of course, the neighborhood $V$
depends on $\tau_1$.) Let $L_{(q_1,t_1)}$ be the length function on $V$
obtained by taking the ${\mathcal L}$-lengths of geodesics parameterized by
points of $\widetilde V$. Let $\varphi_{\tau_1}\colon V\to \Ar$ be defined by
$$\varphi_{\tau_1}(q',\tau')=\frac{1}{2\sqrt{\tau'}}\left({\mathcal
L}(\gamma|_{[0,\tau_1]})+L_{(q_1,t_1)}(q',T-\tau')\right).$$
Clearly, $\varphi_{\tau_1}$ is an upper barrier for $l_x$ at
$(q,\tau)$. According to Lemma~\ref{linequal} we have
\begin{eqnarray*}
\frac{\partial \varphi_{\tau_1}}{d\tau}(q,\tau)+\triangle
\varphi_{\tau_1}(q,\tau) & \le &
\frac{n}{2\sqrt{\tau}(\sqrt{\tau}-\sqrt{\tau_1})}-\frac{\varphi_{\tau_1}(q,\tau)}{\tau}
+\frac{1}{2\tau^{3/2}}{\mathcal L}(\gamma|_{[0,\tau_1]})\\
& & +\frac{K^{\tau}_{\tau_1}(\gamma)}{2\tau^{3/2}}-\frac{{\mathcal
K}_{\tau_1}^{\tau}(\gamma)}{2\sqrt{\tau}(\sqrt{\tau}-\sqrt{\tau_1})^{2}}
\\ & & -\frac{1}{2}\left(\frac{\tau_1}{\tau}\right)^{3/2}\left(R(q_1,t_1)+|X(\tau_1)|^2\right).
\end{eqnarray*}
By Lemma~\ref{Kequal},  it follows easily that
$$\mathrm{lim}_{\tau_1\rightarrow 0^+}
\frac{\partial \varphi_{\tau_1}}{\partial \tau}(q,t)+\triangle
\varphi_{\tau_1}(q,t)\le \frac{(n/2)-l_x(q,t)}{\tau}.$$ This
establishes the result.
\end{proof}

\begin{remark}[Horizontal Laplacian convention and the barrier-inequality limit]
\label{rem:horizontal-laplacian-convention}
The operator $\triangle$ is the horizontal Laplacian on the slice
$M\times\{t=T-\tau\}$ with metric $g(T-\tau)$.
\end{remark}

\begin{claim}[l_x achieves its minimum and the minimizing set is compact]
\label{claim:l-min-achieved-compact}
\uses{def:reduced-length-global}
\label{8.1}
For every $\tau\in(0,T)$, $l_x(\cdot,\tau)$ achieves its minimum. For
every compact interval $I\subset(0,T)$, the set
\[
\{(q,\tau)\in M\times I:l_x(q,\tau)=l_\mathrm{min}(\tau)\}
\]
is compact.
\end{claim}

\begin{proof}\sketch
\uses{cor:l-geodesic-speed-bound}
Metric comparison gives $C^{-1}g(t')\le g(t)\le Cg(t')$. On compact
$I\subset(0,T)$, the min--max estimate gives
$|\sqrt{\tau}X_\gamma(\tau)|\le C_2$ for minimizing curves, hence
$|X_\gamma(\tau)|_{g(T)}\le C_2\sqrt C/\sqrt\tau$ and
\[
\int_0^{\bar\tau}|X_\gamma(\tau)|_{g(T)}d\tau
\le2\sqrt{\bar\tau}C_2\sqrt C.
\]
Thus minimizing endpoints lie in a fixed compact ball $B_T(p,A)$.
Closedness gives compactness of the minimizing set, and continuity of
$l_x$ gives attainment at each fixed $\tau$.
\end{proof}

\begin{theorem}[Upper bound n/2 for the minimum of l_x^tau]
\label{thm:min-l-bound-n-half}
\uses{def:reduced-length-global}
\label{n/2} Suppose that $(M,g(t)),\ 0\le t\le T<\infty$,
is an $n$-dimensional Ricci flow, complete of bounded curvature.
Then for any $x=(p,T)\in M\times [0,T]$ and for every $0<\tau<T$
there is a point $q_\tau\in M$ such that $l_{x}(q_\tau,\tau)\le
\frac{n}{2}$.
\end{theorem}

\begin{proof}
  \uses{claim:l-min-achieved-compact, prop:l-barrier-inequality, prop:forward-difference-maximum, def:reduced-length-global, cor:l-geodesic-speed-bound}\sketch

Set $l_\mathrm{min}(\tau)=\min_{q\in M}l_x(q,\tau)$ using
Claim~\ref{8.1}. Compactness implies continuity of $l_\mathrm{min}$. At a
minimizer $q$ an upper barrier satisfies
\[
\frac{d\varphi}{d\tau}(q,\tau)+\triangle\varphi(q,\tau)
\le\frac{(n/2)-l_\mathrm{min}(\tau)}{\tau}+\epsilon,
\]
while spatial minimality gives $\triangle\varphi(q,\tau)\ge0$. Hence
\[
\limsup_{\tau'\to\tau^+}\frac{l_\mathrm{min}(\tau')-l_\mathrm{min}(\tau)}{\tau'-\tau}
\le\frac{(n/2)-l_\mathrm{min}(\tau)}{\tau}.
\]
Since $l_\mathrm{min}(\tau)\to0$ as $\tau\to0$ (the constant spatial path
has ${\mathcal L}$-length $O(\tau^{3/2})$), the differential inequality
propagates $l_\mathrm{min}(\tau)\le n/2$ to every $\tau<T$.
\end{proof}


\begin{corollary}[Compactness criterion giving the n/2 bound for generalized Ricci flows]
\label{cor:generalized-flow-n-half-bound}
\uses{def:reduced-length-global}
\label{compactcoro}
Suppose that $({\mathcal M},G)$ is a generalized $n$-dimensional Ricci flow and
that $x\in {\mathcal M}$ is given and set $t_0={\bf t}(x)$. We suppose that
there is an open subset $U\subset {\bf t}^{-1}(-\infty,t_0)$ with the following
properties:
\begin{enumerate}
\item For every $y\in U$ there is a minimizing ${\mathcal L}$-geodesic from $x$
to $y$.
\item There are $r>0$ and $\Delta t>0$ such that the backward parabolic neighborhood
$P(x,t_0,r,-\Delta t)$ of $x$ exists in ${\mathcal M}$ and has  the property
that $P\cap {\bf t}^{-1}(-\infty,t_0)$ is contained in $U$.
\item For each compact interval (including the case of degenerate intervals
consisting of a single point) $I\subset (-\infty,t_0)$ the subset of points
$y\in {\bf t}^{-1}(I)\cap U$ for which ${\mathcal L}(y)=\mathrm{inf}_{z\in {\bf
t}^{-1}({\bf t}(y))\cap U}{\mathcal L}(z)$ is compact and non-empty.
\end{enumerate}
 Then
for every $t<t_0$ the minimum of the restriction of $l_x$ to the time-slice
${\bf t}^{-1}(t)\cap U$ is at most $n/2$.
\end{corollary}





\subsection{Extension of the other inequalities in
 Corollary~{\ref{lformula}}}

The material in this subsection is adapted from \cite{Ye}. It
captures (in a weaker way) the fact that, in the case of geodesics
on a Riemannian manifold, the interior of the cut locus in $T_xM$ is
star-shaped from the origin.


\begin{theorem}[Weak (distributional) differential inequalities for l_x]
\label{thm:l-weak-inequalities}
\uses{def:reduced-length-global, cor:reduced-length-pde-inequalities}
\label{weak}
Let $(M,g(t)),\ 0\le t\le T<\infty$, be a Ricci flow, complete and
of bounded curvature, and let $x=(p,T)\in M\times[0,T]$. The last
two inequalities in Corollary~\ref{lformula}, namely
$$\frac{\partial l_x}{\partial \tau}+|\nabla l^\tau_x|^2-R+\frac{n}{2\tau}
- \triangle l^\tau_x\ge 0$$
$$-|\nabla l^\tau_x|^2+R+\frac{l^\tau_x-n}{\tau}+2\triangle l^\tau_x\le 0$$
hold in the weak or distributional sense on all of $M\times\{\tau\}$
for all $\tau>0$. This means that for any $\tau>0$ and for any
non-negative, compactly supported, smooth function $\phi(q)$ on $M$
we have the following two inequalities:
$$\int_{M\times \{\tau\}}\Bigl[ \phi\cdot\left(\frac{\partial l_x}{\partial \tau}+|\nabla
l_x^\tau|^2-R+\frac{n}{2\tau}\right) - l_x^\tau\triangle \phi\Bigr]\
d\vol(g(t))\ge 0$$
$$\int_{M\times \{\tau\}}\Big[ \phi\cdot\left(-|\nabla
l_x^\tau|^2+R+\frac{l_x^\tau-n}{\tau}\right)+2l^\tau_x\triangle
\phi\Bigr]\ d\vol(g(t))\le 0.$$ Furthermore, equality holds in
either of these weak inequalities for all functions $\phi$ as above
and all $\tau$ if and only if it holds in both. In that case $l_x$
is a smooth function on space-time and the equalities hold in the
usual smooth sense.
\end{theorem}
\begin{proof}\sketch
\uses{cor:l-tau-weak-laplacian-inequality,cor:reduced-length-pde-inequalities,prop:l-exp-lipschitz-complete}
Use a partition of unity subordinate to the balls $B'_{(q,t)}$. The local
comparison inequality for $l_x^\tau$ gives the first distributional
inequality, and the second follows from the smooth identity on the full
measure set ${\mathcal U}_x(\tau)$ and the relation $D_2=-2D_1$. Elliptic
regularity and bootstrapping give the final smoothness assertion when
one equality holds for every test function and every $\tau$.
\end{proof}








\begin{remark}[Interpretation of the terms in the weak inequalities for l_x]
\label{rem:l-weak-inequalities-interpretation}
\uses{thm:l-weak-inequalities}

The terms in these inequalities are interpreted in the following
way: First of all, $\nabla l^\tau_x$ and $\triangle l^\tau_x$ are
computed using only the spatial derivatives (i.e., they are
horizontal differential operators). Secondly, since $l_x$ is a
locally Lipschitz functions defined on all of $M\times (0,T)$ we
have seen that $\partial l_x/\partial t$ and $|\nabla l^\tau_x|^2$
are continuous functions on the open subset ${\mathcal U}_x(\tau)$
of full measure in $M\times\{\tau\}$ and furthermore, that they are
locally bounded on all of $M\times\{\tau\}$ in the sense that for
any $q\in M$ there is a neighborhood $V$ of $q$ such that the
restriction of $|\nabla l_x^\tau|^2$ to $V\cap {\mathcal U}_x(\tau)$
is bounded. This means that $\partial l_x/\partial t$ and $|\nabla
l_x^\tau|^2$ are elements of $L_\mathrm{loc}^\infty(M)$ and hence can
be integrated against any smooth function with compact support. In
particular, they are distributions.

 Since
$\nabla l^\tau_x$ is a smooth, locally bounded vector field on an
open subset of full measure, for any compactly supported test
function $\phi$, integration by parts yields
$$\int\triangle \phi\cdot l^\tau_x\ d\vol(g(t))=-\int \langle \nabla \phi,\nabla
l_x^\tau\rangle\ d\vol(g(t)).$$ Thus, formulas in
Theorem~\ref{weak} can also be taken to mean:
$$\int_{M\times \{\tau\}} \Bigl[\phi\cdot(\frac{\partial l_x}{\partial \tau}
+|\nabla l^\tau_x|^2-R+\frac{n}{2\tau})+ \langle \nabla
l_x^\tau,\nabla \phi \rangle\Bigr] d\vol(g(t))\ge 0$$
$$\int_{M\times\{\tau\}} \Bigl[\phi\cdot(-|\nabla
l^\tau_x|^2+R+\frac{l^\tau_x-n}{\tau})-2\langle\nabla
l^\tau_x,\nabla\phi \rangle\Bigr] d\vol(g(t))\le 0.$$
\end{remark}

The rest of this subsection is devoted to the proof of these
inequalities. We fix $(M,g(t)),\ 0\le t\le T<\infty$, as in the
statement of the theorem.
 We
fix $x$ and denote by $L$ and $l$ the functions $L_{x}$ and $l_{x}$.
We also fix $\tau$, and we denote by $L^\tau$ and $l^\tau$ the
restrictions of $L$ and $l$ to the slice $M\times \{T-\tau\}$. We
begin with a lemma.

\begin{lemma}[Uniform gradient and Hessian bound for an upper barrier of L_x]
\label{lem:l-barrier-hessian-bound}
\uses{def:reduced-length-global}
\label{phisupport}
There is a continuous function $C\colon M\times (0,T)\to \Ar$ such
that for each point $(q,t)\in M\times (0,T)$, setting $\tau=T-t$,
the following holds. There is an upper barrier $\varphi_{(q,t)}$ for
$L_x^\tau$ at the point $q$ defined on a neighborhood $U_{(q,t)}$ of
$q$ in $M$ satisfying $|\nabla \varphi_{(q,t)}(q)|\le C(q,t)$ and
$$\Hess(\varphi)(v,v)\le C(q,t)|v|^2$$
for all tangent vectors $v\in T_qM$.
\end{lemma}

\begin{proof}\sketch
\uses{thm:shi-first-derivative,lem:l-geodesic-existence,prop:l-exp-lipschitz-complete,cor:l-length-gradient,cor:l-geodesic-speed-bound}
By Proposition~\ref{lipcomplete}, $L$ is a locally Lipschitz
function on $M\times (0,T)$, and in particular is continuous. The
bound $C(q,t)$ will depend only on the bounds on curvature and its
first two derivatives and on the function $L(q,t)$. Fix $(q,t)$ and
let $\gamma$ be a minimizing ${\mathcal L}$-geodesic from $x$ to
$(q,t)$. (The existence of such a minimizing geodesic is established
in Lemma~\ref{geoexist}.) Fix $\tau_1>0$, with $\tau_1<(T-t)/2$, let
$t_1=T-\tau_1$, and let $q_1=\gamma(\tau_1)$. Consider
$\varphi_{(q,t)}={\mathcal L}(\gamma_{[0,\tau_1]})+
L^\tau_{(q_1,t_1)}$. This is an upper barrier  for $L_x^\tau$ at $q$
defined in some neighborhood $V\subset M$ of $q$.  Clearly, $\nabla
\varphi_{(q,t)}=\nabla L^\tau_{(q_1,t_1)}$ and $\Hess(\varphi_{(q,t)}=\Hess( L^\tau_{(q_1,t_1)})$.



According to Corollary~\ref{DL} we have $\nabla
L^\tau_{(q_1,t_1)}(q)=2\sqrt{\tau}X_{\gamma}(\tau)$. On the other
hand, by Corollary~\ref{minmax} there is a bound on
$\sqrt{\tau}|X_{\gamma}(\tau)|$ depending only on the bounds on
curvature and its first derivatives, on $\tau$ and $\tau_1$ and on
$l_x(q,\tau)$. Of course, by Shi's theorem (Theorem~\ref{firstshi})
for every $\epsilon>0$ the norms of the first derivatives of
curvature on $M\times [\epsilon,T]$ are bounded in terms of
$\epsilon$ and the bounds on curvature.
 This
proves that $|\nabla \varphi_{(q,t)}(q)|$ is bounded by a continuous
function $C(q,t)$ defined on all of $M\times (0,T)$.

 Now consider Inequality~(\ref{Hessin}) for $\gamma$ at $\bar\tau=\tau$.
 It is clear that the first
 two terms on the right-hand side are bounded by $C|Y(\tau)|^2$,
 where $C$ depends on the curvature bound and on $T-t$. We consider
 the last term, $\int_{\tau_1}^{\tau}\sqrt{\tau'}H(X,Y)d\tau'$.
 We claim that this integral is also bounded by $C'|Y(\tau)|^2$
 where $C'$ depends on the bounds on curvature and its first and
 second derivatives along $\gamma_1$ and on $T-t$.
We consider $\tau'\in [\tau_1,\tau]$. Of course,
 $\sqrt{\tau'}|X(\tau')|$ is bounded on this interval. Also,
 $$|Y(\tau')|=\left(\frac{\sqrt{\tau'}-\sqrt{\tau_1}}{\sqrt{\tau}
 -\sqrt{\tau_1}}\right)|Y(\tau)|\le \frac{\sqrt{\tau'}}{\sqrt{\tau}}|Y(\tau)|.$$
Hence $|Y(\tau')|/\sqrt{\tau'}$ and $|Y(\tau')||X(\tau')|$ are
bounded in terms of $T-t$, $|Y(\tau)|$, and the bound on
$\sqrt{\tau'}|X(\tau')|$ along the ${\mathcal L}$-geodesic. From
this it follows immediately from Equation~(\ref{Heqn}) that $H(X,Y)$
is bounded along the ${\mathcal L}$-geodesic  by $C|Y(\bar\tau)|^2$
 where the constant $C$ depends on $T-t$ and the
 bounds on curvature and its first two derivatives.
\end{proof}

Of course, if $(q,t)\in {\mathcal U}_{x}(\tau)$, then this argument
shows that the Hessian of $L^\tau_{x}$ is bounded near $(q,t)$.


At this point in the proof of Theorem~\ref{weak} we wish to employ
arguments using convexity. To carry these out we find it convenient
to work with a Euclidean metric and usual convexity rather than the
given metric $g(t)$ and convexity measured using $g(t)$-geodesics.
In order to switch to a Euclidean metric we must find one that well
approximates $g(t)$. The following is straightforward to prove.

\begin{claim}[Euclidean-approximating ball with controlled Christoffel symbols]
\label{claim:euclidean-approximating-ball}
\uses{lem:l-barrier-hessian-bound}

For each point $(q,t)\in M\times(0,T)$ there is an open metric ball
$B_{(q,t)}$ centered at $q$ in $(M,g(t))$ which is the diffeomorphic
image of a ball $\widetilde B\subset T_qM$ under the exponential map
for $g(t)$ centered at $q$ such that the following hold:
\begin{enumerate}
\item $B_{(q,t)}\subset U_{(q,t)}$ so that the upper barrier $\varphi_{(q,t)}$
from Lemma~\ref{phisupport} is defined on all of $B_{(q,t)}$.
\item The constants $C(z,t)$ of Lemma~\ref{phisupport} satisfy
$C(z,t)\le 2C(q,t)$ for all $z\in B_{(q,t)}$.
\item The push-forward, $h$, under the exponential mapping of the Euclidean
metric on $T_qM$ satisfies $$h/2\le g\le 2h.$$
\item The Christoffel symbols $\Gamma_{ij}^k$ for the
metric $g(t)$ written using the Gaussian normal coordinates (the
image under the exponential mapping of orthonormal linear
coordinates on $T_qM$) are bounded in absolute value by
$1/(8n^3C(q,t))$ where $n$ is the dimension of $M$.
\end{enumerate}
\end{claim}

Instead of working in the given metric $g(t)$ on $B_{(q,t)}$ we
shall use the Euclidean metric $h$ as in the above claim. For any
function $f$ on $B_{(q,t)}$ we denote by $\Hess(f)$ the Hessian
of $f$ with respect to the metric $g(t)$ and by $\Hess^h(f)$
the Hessian of $f$ with respect to the metric $h$. By
Formula~(\ref{Hessian}), for any $z\in B_{(q,t)}$ and any $v\in
T_zM$, we have
$$\Hess(\varphi_{(z,t)})(v,v)=\Hess^h(\varphi_{(z,t)})(v,v)
-\sum_{i,j,k}v^iv^j\Gamma_{ij}^k\frac{\partial
\varphi_{(z,t)}}{\partial x^k}.$$ Thus, it follows from the above
assumptions on the $\Gamma_{ij}^k$ and the bound on $|\nabla
\varphi_{(z,t)}|$ that for all $z\in B_{(q,t)}$ we have
\begin{equation}\label{Hessdiff}
\left|\Hess(\varphi_{(z,t)})(v,v)-\Hess^h(\varphi_{(z,t)})(v,v)\right|\le
\frac{1}{4}|v|^2_h,\end{equation} and hence for every $z\in
B_{(q,t)}$ we have
$$\Hess^h(\varphi_{(z,t)})(v,v)\le 2C(q,t)|v|^2_g+\frac{|v|_h^2}{4}\le
\left(4C(q,t)+\frac{1}{4}\right)|v|_h^2.$$




This means:

\begin{claim}[Upper barrier with strictly negative Euclidean Hessian]
\label{claim:strictly-concave-barrier}
\uses{claim:euclidean-approximating-ball, lem:l-barrier-hessian-bound}

For each $(q,t)\in M\times(0,T)$ there is a smooth function
$$\psi_{(q,t)}\colon B_{(q,t)}\to \Ar$$
 with the property that at each
$z\in B_{(q,t)}$ there is an upper barrier $b_{(z,t)}$ for
$L^\tau+\psi_{(q,t)}$ at $z$ with
$$\Hess^h(b_{(z,t)})(v,v)\le -3|v|_h^2/2$$
for all $v\in T_zM$.
\end{claim}

\begin{proof}\sketch
 Set
$$\psi_{(q,t)}=-(2C(q,t)+1)d^2_h(q,\cdot).$$
Then  for any $z\in B_{(q,t)}$ the function
$b_{(z,t)}=\varphi_{(z,t)}+\psi_{(q,t)}$ is an upper barrier  for
$L^\tau+\psi_{(q,t)}$ at $z$. Clearly, for all $v\in T_zM$ we have
$$\Hess^h(b_{(z,t)})(v,v)=\Hess^h(\varphi_{(z,t)})(v,v)+\Hess^h(\psi_{(q,t)})(v,v)\le -3|v|_h^2/2.$$
\end{proof}

This implies that if $\alpha\colon [a,b]\to B_{(q,t)}$ is any
Euclidean straight-line segment in $B_{(q,t)}$ parameterized by
Euclidean arc length and if $z= \alpha(s)$ for some $s\in (a,b)$,
then
$$(b_{(z,t)}\circ\alpha)''(s)\le -3/2.$$


\begin{claim}[One-dimensional barrier lemma via the maximum principle]
\label{claim:one-dim-barrier-maximum-principle}
Suppose that $\beta\colon [-a,a]\to \Ar$ is a continuous function
and that at each $s\in (-a,a)$ there is an upper barrier $\hat b_s$
for $\beta$  at $s$ with $\hat b_s''\le -3/2$. Then
$$\frac{\beta(a)+\beta(-a)}{2}\le \beta(0)-\frac{3}{4}a^2.$$
\end{claim}

\begin{proof}\sketch
Fix $c<3/4$ and define a continuous function
$$A(s)=\frac{(\beta(-s)+\beta(s))}{2}+cs^2-\beta(0)$$ for $s\in [0,a]$.
Clearly, $A(0)=0$. Also, using the upper barrier at $0$ we see that
for $s>0$ sufficiently small $A(s)<0$. For any $s\in (0,a)$ there is
an upper barrier $c_s=(\hat b_{s}+\hat b_{-s})/2+cs^2-\beta(0)$ for
$A(s)$ at $s$, and $c_s''(t)\le 2c-3/2<0$. By the maximum principle
this implies that $A$ has no local minimum in $(0,a)$, and
consequently that it is a non-increasing function of $s$ on this
interval. That is to say, $A(s)<0$ for all $s\in (0,a)$ and hence
$A(a)\le 0$, i.e., $(\beta(a)+\beta(-a))/2+ca^2\le \beta(0)$. Since
this is true for every $c<3/4$, the result follows.
\end{proof}

Now applying this to Euclidean intervals in $B_{(q,t)}$ we conclude:

\begin{corollary}[Uniform strict convexity of L^tau plus psi]
\label{cor:l-tau-uniform-convexity}\label{lineineq}
\uses{claim:strictly-concave-barrier, claim:one-dim-barrier-maximum-principle}

For any $(q,t)\in M\times (0,T)$, the function
$$\beta_{(q,t)}=L^\tau+\psi_{(q,t)}\colon B_{(q,t)}\to\Ar$$ is
uniformly strictly convex with respect to $h$. In fact, let
$\alpha\colon [a,b]\to B_{(q,t)}$ be a Euclidean geodesic arc. Let
$y,z$ be the endpoints of $\alpha$, let $w$ be its midpoint, and let
$|\alpha|$ denote the length of this arc (all defined using the
Euclidean metric). We have
$$\beta_{(q,t)}(w)\ge \frac{\left(\beta_{(q,t)}(y)+\beta_{(q,t)}(z)\right)}{2}+\frac{3}{16}|\alpha|^2.$$
\end{corollary}

What follows is a simple interpolation result (see \cite{GreeneWu}).
For each $q\in M$ we let $B_{(q,t)}'\subset B_{(q,t)}$ be a smaller
ball centered at $q$, so that $B'_{(q,t)}$ has compact closure in
$B_{(q,t)}$.



\begin{claim}[Smooth strictly convex approximations to beta_(q,t)]
\label{claim:smooth-convex-approximation}
\uses{cor:l-tau-uniform-convexity}

Fix $(q,t)\in M\times(0,T)$, and let $\beta_{(q,t)}\colon B_{(q,t)}\to \Ar$ be
as above. Let $S\subset M$ be the singular locus of $L^\tau$, i.e.,
$S=M\setminus {\mathcal U}_{x}(\tau)$. Set $S_{(q,t)}=B_{(q,t)}\cap S$. Of
course, $\beta_{(q,t)}$ is smooth on $B_{(q,t)}\setminus S_{(q,t)}$. Then there
is a sequence of smooth functions $\{f_k\colon B'_{(q,t)}\to
\Ar\}_{k=1}^\infty$ with the following properties:
\begin{enumerate}
\item As $k\rightarrow \infty$ the functions $f_k$ converge uniformly to $\beta_{(q,t)}$ on
$B_{(q,t)}'$.
\item For any $\epsilon>0$ sufficiently small, let $\nu_\epsilon(S_{(q,t)})$ be the
$\epsilon$-neighborhood (with respect to the Euclidean metric) in
$B_{(q,t)}$ of $S_{(q,t)}\cap B_{(q,t)}$. Then, as $k\rightarrow
\infty$ the restrictions of $f_k$ to $B_{(q,t)}'\setminus
\left(B_{(q,t)}'\cap \nu_\epsilon(S_{(q,t)})\right)$ converge
uniformly in the $C^\infty$-topology to the restriction of
$\beta_{(q,t)}$ to this subset.
\item For each $k$, and for any $z\in B_{(q,t)}'$ and any $v\in T_zM$ we have
$$\Hess(f_k)(v,v)\le -|v|_{g(t)}^2/2.$$
That is to say, $f_k$ is strictly convex with respect to the metric
$g(t)$.
\end{enumerate}
\end{claim}

\begin{proof}\sketch
\uses{cor:l-tau-uniform-convexity}
Fix $\epsilon>0$ sufficiently small so that for any $z\in
B_{(q,t)}'$ the Euclidean $\epsilon$-ball centered at $z$ is
contained in $B_{(q,t)}$. Let $B_0$ be the ball of radius $\epsilon$
centered at the origin in $\Ar^n$ and let $\xi\colon B_0\to \Ar$ be
a non-negative $C^\infty$-function with compact support and with
$\int_{B_0}\xi d\vol_h=1$. We define
$$\beta_{(q,t)}^\epsilon(z)=\int_{B_0}\xi(y)\beta_{(q,t)}(z+y)dy,$$
for all $z\in B_{(q,t)}'$. It is clear that for each $\epsilon>0$
sufficiently small, the function $\beta_{(q,t)}^\epsilon\colon
B_{(q,t)}'\to \Ar$ is $C^\infty$ and that as $\epsilon\rightarrow 0$
the $\beta_{(q,t)}^\epsilon$ converge uniformly on $B_{(q,t)}'$ to
$\beta_{(q,t)}$. It is also clear that for every $\epsilon>0$
sufficiently small, the conclusion of Corollary~\ref{lineineq} holds
for $\beta_{(q,t)}^\epsilon$ and for each Euclidean straight-line
segment $\alpha$ in $B_{(q,t)}'$. This implies that $\Hess^h(\beta_{(q,t)}^\epsilon)(v,v)\le -3|v|_h^2/2$, and hence that
by Inequality~(\ref{Hessdiff}) that
$$\Hess(\beta_{(q,t)}^\epsilon)(v,v)\le -|v|^2_h=-|v|^2_{g(t)}/2.$$
This means that $\beta_{(q,t)}^\epsilon$ is convex with respect to
$g(t)$. Now take a sequence $\epsilon_k\rightarrow 0$ and let
$f_k=\beta_{(q,t)}^{\epsilon_k}$.

Lastly, it is a standard fact that $f_k$ converge uniformly in the
$C^\infty$-topology to $\beta_{(q,t)}$ on any subset of $B_{(q,t)}'$
whose closure is disjoint from $S_{(q,t)}$.
\end{proof}

\begin{definition}[Starred integral over the regular locus]
\label{def:starred-integral}
For any continuous function $\psi$ defined on $B_{(q,t)}'\setminus
\left(S_{(q,t)}\cap B_{(q,t)}'\right)$ we define
$$\int_{(B_{(q,t)}')^*}\psi d\vol(g(t))=\mathrm{lim}_{\epsilon\rightarrow
0}\int_{B_{(q,t)}'\setminus \nu_\epsilon(S_{(q,t)})\cap
B_{(q,t)}'}\psi d\vol(g(t)).$$
\end{definition}

We now have:

\begin{claim}[Weak Laplacian comparison inequality for beta_(q,t)]
\label{claim:beta-weak-laplacian-inequality}
\uses{cor:l-tau-uniform-convexity, def:starred-integral}
\label{wineq}
Let $\phi\colon B_{(q,t)}'\to \Ar$ be a non-negative, smooth
function with compact support. Then
$$\int_{B_{(q,t)}'}\beta_{(q,t)}\triangle \phi d\vol(g(t))\le \int_{(B_{(q,t)}')^*}\phi
\triangle \beta_{(q,t)} d\vol(g(t)).$$
\end{claim}


\begin{proof}\sketch
\uses{claim:smooth-convex-approximation}
Since $f_k\rightarrow \beta_{(q,t)}$ uniformly on $B_{(q,t)}'$ we
have
$$\int_{B_{(q,t)}'}\beta_{(q,t)}\triangle \phi d\vol(g(t))=\mathrm{lim}_{k\rightarrow
\infty}\int_{B_{(q,t)}'}f_k\triangle \phi  d\vol(g(t)).$$ Since
$f_k$ is strictly convex with respect to the metric $g(t)$,
$\triangle f_k\le 0$ on all of $B_{(q,t)}'$. Since $\phi\ge 0$, for
every $\epsilon$ and $k$ we have
$$\int_{\nu_\epsilon(S_{(q,t)})\cap B_{(q,t)}'}\phi\triangle f_k  d\vol(g(t))\le
0.$$ Hence, for every $k$ and for every $\epsilon$ we have
\begin{eqnarray*}
\int_{B_{(q,t)}'}f_k\triangle \phi d\vol(g(t)) & = &
\int_{B_{(q,t)}'}\phi\triangle f_k  d\vol(g(t)) \\
& \le & \int_{B_{(q,t)}'\setminus \left(B_{(q,t)}'\cap
\nu_\epsilon(S_{(q,t)})\right)}\phi\triangle f_k d\vol(g(t)).\end{eqnarray*}
 Taking the limit as $k\rightarrow\infty$,
using the fact that $f_k\rightarrow \beta_{(q,t)}$ uniformly on
$B_{(q,t)}'$ and that restricted to
$B_{(q,t)}'\setminus(B_{(q,t)}'\cap \nu_\epsilon (S_{(q,t)}))$ the
$f_k$ converge uniformly in the $C^\infty$-topology to
$\beta_{(q,t)}$ yields
$$\int_{B_{(q,t)}'}\beta_{(q,t)}\triangle \phi d\vol(g(t))
\le \int_{B_{(q,t)}'\setminus \left(B_{(q,t)}'\cap
\nu_\epsilon(S_{(q,t)})\right)}\phi\triangle \beta_q d\vol(g(t)).$$ Now taking the limit as $\epsilon\rightarrow 0$
establishes the claim.
\end{proof}


\begin{remark}[Laplacian convention with respect to the metric g(t)]
\label{rem:laplacian-convention-metric-gt}
\uses{claim:beta-weak-laplacian-inequality}

Here $\triangle$ denotes the Laplacian with respect to the metric
$g(t)$.
\end{remark}

\begin{corollary}[Weak Laplacian comparison inequality for l^tau]
\label{cor:l-tau-weak-laplacian-inequality}
\uses{def:reduced-length-global, def:starred-integral}
\label{ltauineq}
Let $\phi\colon B_{(q,t)}'\to \Ar$ be a non-negative, smooth
function with compact support. Then
$$\int_{B_{(q,t)}'}l^\tau\triangle \phi d\vol(g(t))\le \int_{(B_{(q,t)}')^*}\phi
\triangle l^\tau d\vol(g(t)).$$
\end{corollary}


\begin{proof}
  \uses{claim:beta-weak-laplacian-inequality, cor:l-tau-weak-laplacian-inequality, cor:reduced-length-pde-inequalities, prop:l-exp-lipschitz-complete}\sketch

Recall that $\beta_{(q,t)}=L^\tau+\psi_{(q,t)}$ and that
$\psi_{(q,t)}$ is a $C^\infty$-function. Hence,
$$\int_{B_{(q,t)}'}\psi_{(q,t)}\triangle \phi d\vol(g(t))= \int_{(B_{(q,t)}')^*}\phi
\triangle \psi_{(q,t)} d\vol(g(t)).$$ Subtracting this equality
from the inequality in the previous claim and dividing by
$2\sqrt{\tau}$ gives the result.
\end{proof}

Now we turn to the proof proper of Theorem~\ref{weak}.

\section{Reduced volume}



We have established that for a Ricci flow $(M,g(t)),\ 0\le t\le T$,
and a point  $x=(p,T)\in M\times [0,T]$ the reduced length function
$l_x$ is defined on all of $M\times (0,T)$. This allows us to
defined the reduced volume of $M\times \{\tau\}$ for any $\tau\in
(0,T)$ Recall that the {\sl reduced volume} of
$M$ is defined to be
$$\widetilde V_x(M,\tau) = \int_{M} \tau^{-\frac{n}{2}}exp(-l_x(q,\tau))dq.$$
This function is defined for $0<\tau<T$.

There is one simple case where we can make an explicit computation.
\begin{lemma}[Reduced volume of flat Euclidean space equals (4pi)^(n/2)]
\label{lem:reduced-volume-flat-space}
\uses{def:reduced-length-global}
\label{flat}
If $(M,g(t))$ is flat Euclidean $n$-space (independent of $t$), then
for any $x\in \Ar^n\times (-\infty,\infty)$ we have
$$\widetilde V_x(M,\tau)=(4\pi)^{n/2}$$ for all $\tau>0$.
\end{lemma}

\begin{proof}\sketch
By symmetry we can assume that $x=(0,T)\in \Ar^n\times [0,T]$, where
$0\in \Ar^n$ is the origin. We have already seen that the ${\mathcal
L}$-geodesics in flat space are the usual geodesics when
parameterized by $s=\sqrt{\tau}$. Thus, for any $X\in
\Ar^n=T_0\Ar^n$ $\gamma_X(\tau)=2\sqrt{\tau}X$, and hence ${\mathcal
L}\mathrm{exp}(X,\bar\tau)=2\sqrt{\bar\tau}X$. This means that for any
$\tau>0$ and any $X\in T_0\Ar^n$ we have ${\mathcal U}(\tau)=T_pM$,
and ${\mathcal J}(X,\tau)=2^n\tau^{n/2}$. Also,
 $L_x(X,\tau)=2\sqrt{\tau}|X|^2$, so that $l_x(X,\tau)=|X|^2$.
 Thus, for any $\tau>0$
 $$\widetilde
 V_x(\Ar^n,\tau)=\int_{\Ar^n}\tau^{-n/2}e^{-|X|^2}2^n\tau^{n/2}dX=(4\pi)^{n/2}.$$
\end{proof}


In the case when $M$ is non-compact, it is not clear {\em a priori}
that the integral defining the reduced volume is finite in general.
In fact, as the next proposition shows, it is always finite and
indeed, it is bounded above by the integral for  $\Ar^n$.

\begin{theorem}[Reduced volume is finite, bounded by (4pi)^(n/2), and non-increasing]
\label{thm:reduced-volume-monotone-bound}
\uses{def:reduced-length-global}
\label{4pi} Let $(M,g(t)),\ 0\le t\le T$, be a Ricci flow
of bounded curvature with the property that for each $t\in [0,T]$
the Riemannian manifold $(M,g(t))$ is complete. Fix a point
$x=(p,T)\in M\times [0,T]$. For every $0<\tau<T$ the reduced volume
$$\widetilde V_x(M,\tau)=\int_{M} \tau^{-\frac{n}{2}}\mathrm{exp}(-l_x(q,\tau))dq$$
is absolutely convergent and  $\widetilde V_x(M,\tau)\le
(4\pi)^{\frac{n}{2}}$. The function $\widetilde V_x(M,\tau)$ is a
non-increasing function of $\tau$ with
$$\mathrm{lim}_{\tau\rightarrow 0}\widetilde V_x(M,\tau)=(4\pi)^{\frac{n}{2}}.$$
\end{theorem}

\begin{proof}\sketch
\uses{prop:l-exp-lipschitz-complete,lem:reduced-volume-pullback,prop:reduced-volume-integrand-monotone,thm:reduced-volume-monotone,cor:injectivity-set-contains-ball-small-tau}
By Proposition~\ref{lipcomplete} ${\mathcal U}_x(\tau)$ is an open
subset of full measure in $M$. Hence, $$\widetilde
V_x(M,\tau)=\int_{{\mathcal U}_x(\tau)}\tau^{-\frac{n}{2}}\mathrm{exp}(-l_x(q,\tau))dq.$$
 Take
linear orthonormal coordinates $(z^1,\ldots,z^n)$ on $T_pM$. It
follows from  the previous equality and Lemma~\ref{redvoltilde} that
$$\widetilde V_x(M,\tau)=\int_{\widetilde {\mathcal U}_{x}(\tau)}f(Z,\tau)dz^1\cdots dz^n,$$
where $f(Z,\tau)=\tau^{-\frac{n}{2}}e^{-\widetilde
l(Z,\tau)}{\mathcal J}(Z,\tau)$. By Proposition~\ref{fmonotone} for
each $Z$ the integrand, $f(Z,\tau)$, is a non-increasing function of
$\tau$ and the function converges uniformly on compact sets as
$\tau\rightarrow 0$ to $2^ne^{-|Z|^2}$. This implies that
$f(Z,\tau)\le 2^ne^{-|Z|^2}$ for all $\tau>0$, and hence that
$$\int_{\widetilde{\mathcal U}_x(\tau)}f(Z,\tau)dz^1\ldots dz^n$$
converges absolutely for each $\tau>0$, and the integral has value
at most $(4\pi)^{n/2}$.

Fix $0<\tau_0<T$. According to Theorem~\ref{Amono} (with $A=M\times
(T-\tau_0,T)$), the reduced volume $\widetilde V_x(M,\tau)$ is a
non-increasing function of $\tau$ on $(0,\tau_0]$. Since this is
true for any $0<\tau_0<T$, it follows that $\widetilde V_x(M,\tau)$
is a non-increasing function of $\tau$ for all $\tau\in (0,T)$.
 (This of course is a consequence of the
monotonicity of $f(Z,\tau)$ in $\tau$ and the fact that $\widetilde {\mathcal
U}_x(\tau)\subset \widetilde {\mathcal U}_x(\tau')$ for $\tau'<\tau$.)



To show that $\mathrm{lim}_{\tau\rightarrow 0}\widetilde
 V_x(M,\tau)= (4\pi)^{n/2}$ we need only see that for
  each $A<\infty$ for all $\tau>0$ sufficiently small
 $\widetilde {\mathcal U}_x(\tau)$ contains the ball of radius $A$ centered at the origin
 in $T_pM$. Since the  curvature is bounded, this is exactly
 the content of Corollary~\ref{Jlimit1}.
\end{proof}



\subsection{Converse to Lemma~{\ref{flat}}}

In Lemma~\ref{flat} we showed that for the trivial flow on flat
Euclidean $n$-space and for any point $x\in \Ar^n\times \{T\}$ the
reduced volume $\widetilde V_x(\Ar^n,\tau)$ is independent of
$\tau>0$  and is equal to $(4\pi)^{n/2}$. In this subsection we use
the monotonicity results of the last subsection to establish the
converse to Lemma~\ref{flat}, namely to show that if $(M,g(t)),0\le
t\le T$, is a Ricci flow complete with bounded curvature and if
$\widetilde V_x(M,\bar\tau)=(4\pi)^{n/2}$ for some $\bar\tau>0$ and
some $x\in M\times\{T\}$, then the flow on the interval
$[T-\bar\tau,T]$ is the trivial flow on flat Euclidean $n$-space.



\begin{proposition}[Reduced volume equal to (4pi)^(n/2) forces the trivial flat flow]
\label{prop:reduced-volume-rigidity}\label{4picase}
\uses{def:reduced-length-global}

Suppose that $(M,g(\tau)),\ 0\le \tau\le T$,
 is a solution to the backward Ricci flow
equation, complete of bounded curvature. Let $x=(p,T)\in M\times
\{T\}$, and suppose that $0<\bar\tau<T$. If $\widetilde
V_x(M,\bar\tau)=(4\pi)^{n/2}$, then the backward Ricci flow on the
interval $[0,\bar\tau]$ is the trivial flow on flat Euclidean space.
\end{proposition}

\begin{proof}\sketch
\uses{thm:reduced-volume-monotone-bound,lem:reduced-volume-flat-space,prop:reduced-volume-integrand-monotone,prop:l-length-hessian-inequality,prop:l-length-laplacian-inequality}
If $\widetilde V_x(M,\bar\tau)=(4\pi)^{n/2}$, then by
Lemma~\ref{flat}, $\widetilde V_x(M,\tau)$ is constant on the
interval $(0,\bar\tau]$. Hence, it follows from the proof of
Theorem~\ref{4pi} that the closure of $\widetilde{\mathcal U}(\tau)$
is all of $T_pM$ for all $\tau\in (0,\bar\tau]$ and that
$f(Z,\tau)=e^{-|Z|^2}2^n$ for all $Z\in T_pM$ and all $\tau\le
\bar\tau$. In particular,
$$\frac{\partial \mathrm{ln}(f(Z,\tau))}{\partial \tau}=0.$$
From the proof of
 Proposition~\ref{fmonotone} this means that
Inequality~(\ref{Jinequ}) is an equality and consequently, so is
Inequality~(\ref{basic}). Thus, by Proposition~\ref{Hessineq}
 (with $\tau_1=0$) each of the vector fields $Y_\alpha(\tau)=\tilde
Y_\alpha(\tau)$ is both a Jacobi field and adapted. By
Proposition~\ref{triL} we then have
$$\Ric+\Hess(l^\tau_x)=\frac{g}{2\tau}.$$
In particular, $l_x$ is smooth. Let $\varphi_\tau\colon M\to M,\
0<\tau\le \bar\tau$,\ be the one-parameter family of diffeomorphisms
obtained by solving
$$\frac{d\varphi_\tau}{d\tau}=\nabla
l_x(\cdot,\tau)\ \ \ \ \mathrm{and} \ \ \ \varphi_{\bar\tau}=\Id.$$ We now consider
$$h(\tau)=\frac{\bar\tau}{\tau}\varphi_\tau^*g(\tau).$$
We compute \begin{eqnarray*}\frac{\partial h}{\partial \tau} & = &
-\frac{\bar\tau}{\tau^2}\varphi_\tau^*g(\tau)
+\frac{\bar\tau}{\tau}\varphi_\tau^*{\mathcal
L}_{\frac{d\varphi_\tau}{d\tau}}(g(\tau))
+\frac{\bar\tau}{\tau}\varphi_t^*2\Ric(g(\tau)) \\
& = & -\frac{\bar\tau}{\tau^2}\varphi_\tau^*g(\tau) +
\frac{\bar\tau}{\tau}\varphi_\tau^*2\Hess(l_x^\tau)+\frac{\bar\tau}{\tau}
\varphi_\tau^*\left(\frac{1}{\tau}g(\tau)-2\Hess(l_x^\tau)\right)=0.\end{eqnarray*} That is to say the family
of metrics $h(\tau)$ is  constant in $\tau$: for all $\tau\in
(0,\bar\tau]$ we have $h(\tau)=h(\bar\tau)=g(\bar\tau)$. It then
follows that
$$g(\tau)=\frac{\tau}{\bar\tau}(\varphi_\tau^{-1})^*g(\bar\tau),$$
which means that the entire flow in the interval $(0,\bar\tau]$
differs by diffeomorphism and scaling from $g(\bar\tau)$. Suppose
that $g(\bar\tau)$ is not flat, i.e., suppose that there is some
$(x,\bar\tau)$ with $|\Rm(x,\bar\tau)|=K>0$. Then from the flow
equation we see that $|\Rm(\varphi^{-1}_\tau(x),\tau)|=K\bar\tau^2/\tau^2$, and these
curvatures are not bounded as $\tau\rightarrow 0$. This is a
contradiction. We conclude that $g(\bar\tau)$ is flat, and hence,
again by the flow equation so are all the $g(\tau)$ for $0<\tau\le
\bar\tau$, and by continuity, so is $g(0)$. Thus, $(M,g(\tau))$ is
isometric to a quotient of $\Ar^n$ by a free, properly discontinuous
group action. Lastly, since $\widetilde V_x(M,\tau)=(4\pi)^{n/2}$,
it follows that $(M,g(\tau))$ is isometric to $\Ar^n$ for every
$\tau\in [0,\bar\tau]$. Of course, it then follows that the flow is
the constant flow.
\end{proof}
